% Part: intuitionistic-logic
% Chapter: tableaux
% Section: rules

\documentclass[../../../include/open-logic-section]{subfiles}

\begin{document}

\olfileid{int}{tab}{rul}

\olsection{قواعد المنطق الحدسي}

قواعد الرابطين~$\land$ و$\lor$ هي نفسها قواعد !!{tableau} القضايا
المعلّمة العادية، مع إضافة السوابق فحسب. ففي كل حالة تنتج القاعدة
المطبقة على !!{formula} معلّمة
$\sFmla{S}{!A}[\sigma]$ !!{formula} جديدة تحمل هي أيضًا السابقة~$\sigma$.
وينبغي أن يكون هذا واضحًا حدسيًا: فمثلًا، إذا كانت~$!A \land !B$
صادقة في (عالم تسمّيه)~$\sigma$، فإن~$!A$ و$!B$ صادقتان
في~$\sigma$ (لا في عالم آخر). نجمع قواعد~$\land$ و$\lor$
في~\olref{tab:prop-rules}.

\begin{table}
  \[\def\arraystretch{3}\begin{array}{|c|c|}
    \hline
    \AxiomC{\sFmla{\True}{!A \land !B}[\sigma]}
    \RightLabel{\TRule{\True}{\land}}
    \UnaryInfC{\sFmla{\True}{!A}[\sigma]}
    \noLine
    \UnaryInfC{\sFmla{\True}{!B}[\sigma]}
    \DisplayProof
    &
    \AxiomC{\sFmla{\False}{!A \land !B}[\sigma]}
    \RightLabel{\TRule{\False}{\land}}
    \UnaryInfC{$\sFmla{\False}{!A}[\sigma] \quad \mid \quad
      \sFmla{\False}{!B}[\sigma]$}
    \DisplayProof
    \\[2ex]
    \hline
    \AxiomC{\sFmla{\True}{!A \lor !B}[\sigma]}
    \RightLabel{\TRule{\True}{\lor}}
    \UnaryInfC{$\sFmla{\True}{!A}[\sigma] \quad \mid \quad
      \sFmla{\True}{!B}[\sigma]$}
    \DisplayProof
    &
    \AxiomC{\sFmla{\False}{!A \lor !B}[\sigma]}
    \RightLabel{\TRule{\False}{\lor}}
    \UnaryInfC{\sFmla{\False}{!A}[\sigma]}
    \noLine
    \UnaryInfC{\sFmla{\False}{!B}[\sigma]}
    \DisplayProof
    \\[2ex]
    \hline
  \end{array}\]
  \caption{قواعد !!{tableau} ذات السوابق للرابطين~$\land$ و$\lor$}
  \ollabel{tab:prop-rules}
\end{table}

شرط الإغلاق شبيه بشرط !!{tableau} العادية، إلا أننا نشترط تطابق السوابق
أيضًا، لا تطابق !!{formula} وحدها. بل يمكننا في الواقع أن نكون أكثر
تساهلًا قليلًا: بما أن !!{formula} تظل صادقة متى صدقت في النماذج
الحدسية، يستحيل أن تكون~$!A$ صادقة عند~$\sigma$ وكاذبة عند أي سابقة
متاحة~$\sigma.{*}$. ولذلك تكون الشعبة مغلقة إذا احتوت، لبعض
السابقة~$\sigma$ و!!{formula}~$!A$، كلًا من
\[
\sFmla{\True}{!A}[\sigma] \quad\text{و}\quad \sFmla{\False}{!A}[\sigma.{*}]
\]
لاحظ أنه إذا عُكست العلامتان، أي إذا احتوت الشعبة
\[
\sFmla{\False}{!A}[\sigma] \quad\text{و}\quad \sFmla{\True}{!A}[\sigma.{*}]
\]
فلا تكون مغلقة إلا إذا كان~$*$ المتتالية الخالية.

وفوق ذلك، تكون الشعبة مغلقة إذا احتوت~$\sFmla{\True}{\bot}[\sigma]$.

وقواعد وضع الافتراضات هي أيضًا كما في !!{tableau} العادية، إلا أننا
نستعمل السابقة~$1$ دائمًا للافتراضات. (ولا يهم أي سابقة نستعمل ما دامت
واحدة لجميع الافتراضات.) لذلك نقول، مثلًا، إن
\[
!B_1, \dots, !B_n \Proves !A
\]
إذا وفقط إذا وُجد جدول مغلق للافتراضات
\[
\sFmla{\True}{!B_1}[1], \dots, \sFmla{\True}{!B_n}[1],
\sFmla{\False}{!A}[1].
\]

أما قواعد الشرط~$\lif$ فتختلف عن نظيراتها في الحالتين الكلاسيكية
والموجهاتية. تمدّد القاعدة~$\TRule{\True}{\lif}$ شعبة تحتوي
$\sFmla{\True}{!A \lif !B}[\sigma]$ بإضافة
$\sFmla{\False}{!A}[\sigma.{*}]$ و$\sFmla{\True}{!B}[\sigma.{*}]$
على شعبتين مختلفتين. ولا يمكن تطبيقها إلا على سابقة~$\sigma.{*}$
\emph{ظهرت من قبل} على الشعبة التي تُطبّق فيها. ولنسمّ سابقة كهذه
«مستعملة» (على الشعبة). (ولأن~$\sigma.{*}$ تشمل~$\sigma$ نفسها،
يمكن دائمًا تطبيق القاعدة بإضافة الصيغتين المعلّمتين ذواتي السابقتين
$\sFmla{\False}{!A}[\sigma]$ و$\sFmla{\True}{!B}[\sigma]$
على شعبتين منفصلتين.)

تمدّد القاعدة~$\TRule{\False}{\lif}$ شعبة تحتوي
$\sFmla{\False}{!A \lif !B}[\sigma]$ بإضافة كل من
$\sFmla{\True}{!A}[\sigma.n]$ و$\sFmla{\False}{!B}[\sigma.n]$
إلى الشعبة نفسها، حيث تكون~$\sigma.n$ سابقة جديدة على الشعبة.

وتُعرّف قواعد~$\lnot$ على نحو مماثل (باستعمال تعريف
$\lnot !A$ بوصفها~$!A \lif \lfalse$).

ترد القواعد في~\olref{tab:rules-lif-lnot}.

\begin{table}
  \[\def\arraystretch{3}\begin{array}{|c|c|}
    \hline
    \AxiomC{\sFmla{\True}{\lnot !A}[\sigma]}
    \RightLabel{\TRule{\True}{\lnot}}
    \UnaryInfC{$\sFmla{\False}{!A}[\sigma.{*}]$}
    \DisplayProof
    &
    \AxiomC{\sFmla{\False}{\lnot !A}[\sigma]}
    \RightLabel{\TRule{\False}{\lnot}}
    \UnaryInfC{\sFmla{\True}{!A}[\sigma.n]}
    \DisplayProof
    \\[1ex]
    \text{$\sigma.{*}$ مستعملة} & \text{$\sigma.n$ جديدة}\\
    \hline
    \AxiomC{\sFmla{\True}{!A \lif !B}[\sigma]}
    \RightLabel{\TRule{\True}{\lif}}
    \UnaryInfC{$\sFmla{\False}{!A}[\sigma.{*}] \quad \mid \quad
      \sFmla{\True}{!B}[\sigma.{*}]$}
    \DisplayProof
    &
    \AxiomC{\sFmla{\False}{!A \lif !B}[\sigma]}
    \RightLabel{\TRule{\False}{\lif}}
    \UnaryInfC{\sFmla{\True}{!A}[\sigma.n]}
    \noLine
    \UnaryInfC{\sFmla{\False}{!B}[\sigma.n]}
    \DisplayProof
    \\[1ex]
    \text{$\sigma.{*}$ مستعملة} & \text{$\sigma.n$ جديدة}\\
    \hline
  \end{array}\]
  \caption{قواعد !!{tableau} ذات السوابق للرابطين~$\lnot$ و$\lif$}
  \ollabel{tab:rules-lif-lnot}
\end{table}

\end{document}
